\section{Technologie}

\subsection{Architektura}

System jest aplikacją internetową złożoną z trzech części: klienta działającego w przeglądarce, serwera udostępniającego REST API oraz bazy danych. Serwer to \textbf{modularny monolit} -- jedna aplikacja podzielona na moduły domenowe, które korzystają ze siebie wyłącznie przez publiczne funkcje serwisów:

\begin{table}[H]
\centering
\begin{tabular}{|l|p{10cm}|}
\hline
\textbf{Moduł} & \textbf{Odpowiedzialność} \\ \hline
\texttt{identity} & konta, role, logowanie, sesja, administracja kont \\ \hline
\texttt{curriculum} & kursy, lekcje, publikacja, klasy i widoczność kursów \\ \hline
\texttt{content} & bloki lekcji (tekst, załączniki z formami V/A/K), pliki \\ \hline
\texttt{learning\_style} & ankieta VAK, profil ucznia, personalizacja lekcji, adaptacja wag \\ \hline
\texttt{assessment} & quiz lekcji, ocenianie odpowiedzi, edytor quizu \\ \hline
\end{tabular}
\caption{Moduły serwera}
\end{table}

Reguły personalizacji, adaptacji i oceniania są czystymi funkcjami, niezależnymi od frameworka i bazy danych. Wynik quizu trafia do modelu ucznia tylko przez jedną funkcję modułu \texttt{learning\_style}, a klient nie ma możliwości samodzielnego zgłoszenia wyniku.

\subsection{Technologie oprogramowania}

\textbf{Klient:}
\begin{itemize}[leftmargin=*]
    \item \textbf{Next.js 16} i \textbf{React 19} -- aplikacja internetowa z routingiem opartym na katalogach
    \item \textbf{TypeScript} -- typowany JavaScript
    \item \textbf{Tailwind CSS 4} i \textbf{shadcn/ui} (Base UI) -- style i komponenty interfejsu
    \item \textbf{TanStack Query} -- pobieranie i buforowanie danych z serwera
    \item \textbf{TipTap} -- edytor tekstu lekcji (formatowanie, tabele, wklejanie z Worda)
    \item \textbf{KaTeX} -- wyświetlanie wzorów matematycznych
    \item \textbf{React Hook Form}, \textbf{react-dropzone}, \textbf{Lucide} -- formularze, wgrywanie plików, ikony
\end{itemize}

\textbf{Serwer:}
\begin{itemize}[leftmargin=*]
    \item \textbf{Python 3.12} i \textbf{FastAPI} -- REST API
    \item \textbf{SQLAlchemy 2} i \textbf{psycopg 3} -- dostęp do bazy danych
    \item \textbf{Alembic} -- migracje schematu bazy
    \item \textbf{Pydantic} -- walidacja danych wejściowych i wyjściowych
    \item \textbf{PyJWT} i \textbf{bcrypt} -- sesja i haszowanie haseł
\end{itemize}

\textbf{Baza danych:}
\begin{itemize}[leftmargin=*]
    \item \textbf{PostgreSQL 17} uruchamiany w kontenerze \textbf{Docker} (Docker Compose); kolumny \textbf{JSONB} dla treści lekcji i quizów
    \item Pliki wgrane przez nauczycieli przechowywane na dysku serwera, poza repozytorium
\end{itemize}

\textbf{Narzędzia deweloperskie:}
\begin{itemize}[leftmargin=*]
    \item \textbf{Git} i \textbf{GitHub} -- kontrola wersji
    \item \textbf{pytest} -- testy serwera; \textbf{Ruff} -- lint i formatowanie kodu Python
    \item \textbf{ESLint} i \textbf{Prettier} -- lint i formatowanie kodu klienta
    \item \textbf{npm} i \textbf{pip} -- menedżery pakietów
\end{itemize}

\subsection{Bezpieczeństwo}

\begin{itemize}[leftmargin=*]
    \item Sesja jako token JWT w ciasteczku \texttt{HttpOnly} (niedostępnym z JavaScriptu), \texttt{SameSite=Lax}
    \item Rola i aktywność konta odczytywane z bazy przy każdym żądaniu -- blokada konta działa natychmiast
    \item Hasła haszowane algorytmem bcrypt
    \item Uprawnienia sprawdzane na serwerze; zasób niedostępny dla użytkownika zwraca 404, bez zdradzania jego istnienia
    \item Wgrywane pliki: biała lista rozszerzeń (bez SVG i HTML), limit 50 MB, losowe nazwy na dysku
    \item Ankieta i quiz wysyłane do klienta bez klucza odpowiedzi; punkty i wynik liczy serwer
    \item CORS tylko dla adresu klienta, sekrety wyłącznie w plikach \texttt{.env} poza repozytorium
\end{itemize}

\subsection{Licencje}

Wszystkie wykorzystane technologie są dostępne na licencjach otwartych, pozwalających na swobodne użycie:

\begin{itemize}[leftmargin=*]
    \item Next.js, React, Tailwind CSS, shadcn/ui, TanStack Query, TipTap, KaTeX -- MIT License
    \item FastAPI, SQLAlchemy, Alembic, Pydantic, pytest -- MIT License
    \item PostgreSQL -- PostgreSQL License (podobna do MIT)
\end{itemize}